\documentclass[12pt,a4paper]{article}

\usepackage[utf8]{inputenc}
\usepackage[T1]{fontenc}
\usepackage{amsmath,amssymb,amsthm,mathtools}
\usepackage[colorlinks=true,linkcolor=blue,citecolor=red,urlcolor=blue]{hyperref}
\renewcommand*{\HyperDestNameFilter}[1]{\jobname-#1}
\usepackage{geometry}
\geometry{margin=2.5cm}
\emergencystretch=1.5em
\usepackage{enumitem}
\usepackage{booktabs}
\usepackage{graphicx}
\usepackage{needspace}
\usepackage[capitalise,noabbrev]{cleveref}

\newtheorem{theorem}{Theorem}[section]
\newtheorem{lemma}[theorem]{Lemma}
\newtheorem{proposition}[theorem]{Proposition}
\newtheorem{corollary}[theorem]{Corollary}
\newtheorem{conjecture}[theorem]{Conjecture}
\theoremstyle{definition}
\newtheorem{definition}[theorem]{Definition}
\theoremstyle{remark}
\newtheorem{remark}[theorem]{Remark}

\title{The Riemann Hypothesis: Cross-Route Synthesis\\[4pt]
\large Part~IV of a Five-Part Series}
\author{Lukas Geiger\\
\small Bernau, Germany\\
\small ORCID: \href{https://orcid.org/0009-0005-7296-1534}{0009-0005-7296-1534}}
\date{\today}

\begin{document}
\overfullrule=0pt
\maketitle

\begin{abstract}\footnotesize
This paper, Part~IV of a five-part series
\cite{Geiger2026partI,Geiger2026partII,Geiger2026partIII,Geiger2026partV},
collects three classes of material that do not fit the scope of Parts~II
(Even Dominance) or III (Spectral Routes): the cross-route synthesis
comparing all three active branches~A, B, Z; the dormant routes and
external lines of the wider landscape; and a self-correcting record of
hypothesis corrections during the iterative research process.

The cross-route synthesis identifies a shared Galerkin pattern across
branches --- each conditional reduction involves a fixed-parameter
positive structure that must be promoted to a uniform-limit statement via
a growing parameter $N(\lambda)$, $d(\lambda)$, or $m(\lambda)$ --- and
tabulates the characteristic constants ($C^* = 64\pi^2/3$ as the
current Branch~A boundary-moment benchmark for $\lambda=5,7$,
$|\kappa_\infty|\lesssim 10^{-4}$ for Branch~B,
$g_0(m) \approx 2\pi/\log\gamma_m$ for Branch~Z).

\medskip\noindent
\textbf{Series status (v3.0, 2026-05-26; criticality-lens update 2026-05-31).}
The series is restructured from a trilogy to a five-part series; the
conditional reduction status from v2.2 is unchanged. A 2026 working phase adds
\cref{sec:sharpened-wall}: a \emph{heuristic} criticality lens
($\mathrm{RH}\Leftrightarrow\Lambda=0$) that strategically re-weights the routes,
a sharper localization of the common wall (the analytic continuation across
$\Re s=1$), a computable diagnostic built from the \emph{classical}
Laguerre--P\'olya/Herglotz criterion, and four honest closures of
\emph{concrete constructions} (not routes, and several of them self-corrections
of internally proposed constructions). None of these closes RH, refutes a route,
or changes the conditional-reduction status.
\end{abstract}

{\footnotesize
\noindent\textbf{MSC 2020:} 11M26, 11M36, 47A75\\
\textbf{Keywords:} Riemann Hypothesis, cross-route synthesis,
Galerkin pattern, dormant routes, hypothesis corrections,
landscape survey
}

\clearpage
\tableofcontents
\clearpage


% =============================================================================
\section{Introduction}
\label{sec:intro}
% =============================================================================

This paper, Part~IV of a five-part series, serves as the ``miscellaneous
and combined'' component of the landscape atlas:
\begin{description}[font=\bfseries,leftmargin=2em]
  \item[Part~I \cite{Geiger2026partI}:] Foundations and Obstructions.
  \item[Part~II \cite{Geiger2026partII}:] Even Dominance (Branch~A).
  \item[Part~III \cite{Geiger2026partIII}:] Spectral Routes
    (Branches~B and~Z).
  \item[Part~IV (this paper):] Cross-Route Synthesis, dormant routes,
    external lines, methodical dead ends, hypothesis corrections.
  \item[Part~V \cite{Geiger2026partV}:] Conclusio.
\end{description}

The three active branches~A (Even Dominance), B (Hadamard Strip / I-1
Normal-Family), and Z (CCM Microcluster) are developed in detail in
Parts~II and~III. This paper identifies the structural patterns they
share, surveys the wider landscape of dormant and external routes, and
provides a transparent record of hypothesis corrections during the
iterative development.


% =============================================================================
\section{Cross-Route Synthesis}
\label{sec:cross-route}
% =============================================================================

\subsection{Characteristic Constants}

Each active branch produces a characteristic constant or scaling law
whose identification anchors the conditional reduction:

\begin{center}
\renewcommand{\arraystretch}{1.3}
\resizebox{\textwidth}{!}{%
\begin{tabular}{l|l|l|l}
\textbf{Branch} & \textbf{Constant} & \textbf{Scaling law} & \textbf{Theoretical identification} \\\hline
A (Even Dom.) & $C^* = 64\pi^2/3 \approx 210.55$ & $\mathrm{sLI} \sim C^*\,\lambda/(NL)$ for the $\lambda=5,7$ benchmark; $\lambda=3$ separate and $\lambda=11$ open & $16\int_0^\infty t^2/\sinh^2(t/2)\,dt$ (reverse-fit candidate; exact collar kernel known) \\
B (Hadamard) & $|\kappa_\infty(11)|\lesssim 10^{-4}$ & corrected (transient conflation); sign open & finite-part residual curvature of Foundation-Lock \\
Z (CCM) & $g_0(m) \approx 2\pi/\log\gamma_m$ & log-decreasing & Montgomery--Odlyzko (classical) \\
\end{tabular}}
\end{center}

\begin{remark}[Identification status]
The constant $C^* = 64\pi^2/3$ for Branch~A is an identification
candidate for the $\lambda=5,7$ benchmark. The 27~May audit shows that the
simple endpoint-kernel factorisation is false; the exact collar kernel is
$K(t)=e^t\cosh(t)/(2\sinh^2t)$, so a rigorous derivation must pass through
the eigenmode asymptotics. The $\lambda=3$ sweep is recorded as a
low-bandwidth outlier, while the single $\lambda=11,N=200$ point
($C_{\mathrm{eff}}\approx168$) remains an open discriminator
(Part~II~\cite{Geiger2026partII}, Remark on derivation candidate). The
constant $\kappa_\infty$ for Branch~B is established empirically in the
conservative range $[1.6,1.8]\times 10^{-3}$ (Part~III~\cite{Geiger2026partIII}). The gap
$g_0(m)$ for Branch~Z rests on classical pair-correlation heuristics.
\end{remark}


\subsection{The Shared Galerkin Pattern}

Despite their different analytical mechanisms, all three branches share a
common structural pattern:

\begin{enumerate}
  \item \textbf{Fixed-parameter positive structure:} At each finite
    $(\lambda, N)$ (or $(\lambda, m)$ for Branch~Z), the relevant
    quantity is computable and has the correct sign/bound.
  \item \textbf{Uniform-limit requirement:} The conditional reduction
    requires promoting the fixed-parameter statement to a uniform
    assertion as $\lambda \to \infty$ with a parameter
    ($N(\lambda)$, $d(\lambda)$, $m(\lambda)$) growing alongside.
  \item \textbf{Transient vs.\ converged regime:} The warm-up parameter is
    branch-dependent. For the GF1/$C^*$ branch the current diagnostic scale is
    $\xi=NL/\lambda$, not $N/L$: the $\lambda=5,7$ rows enter the benchmark
    window near $\xi\approx 50$, whereas the $\lambda=11,N=200$ point remains
    an open pre-convergence discriminator. Other Galerkin branches still track
    their own resolution ratios.
\end{enumerate}

This pattern is consistent with a \emph{Bernstein--Galerkin conjecture}:
the Galerkin truncation of the Weil quadratic form at dimension $N$
converges to the infinite-dimensional form in the relevant operator
topology, and the convergence is uniform in $\lambda$ after a
$\lambda$-dependent warm-up phase.

\begin{remark}[Status of the conjecture]
The Bernstein--Galerkin pattern is an empirical observation from
5--7 data points per branch. It has not been formalised as a
mathematical statement. Its role is heuristic: it provides a
mechanistic explanation for why all three branches have similar
structural gaps (fixed-parameter $\checkmark$, uniform-limit open).
\end{remark}


\subsection{Two Isolated Proof Slots}

The cross-route analysis identifies two concrete, isolated proof slots
that, if resolved, would close one or more branches:

\begin{enumerate}
  \item \textbf{B-1: Uniform log-curvature control from Foundation-Lock.}
    If $H_\lambda(\delta) \leq K\delta^2$ with
    $K \leq 1.8\times 10^{-3}$ is derived from the Foundation-Lock
    structure, then Theorem~B-1
    (Part~III~\cite{Geiger2026partIII}) gives
    $A_\lambda \leq 1.0063$, closing Branch~B.
  \item \textbf{GF1: Eigenmode asymptotics + factor-16 identification.}
    If the exact collar kernel and the $q_a(k)$ eigenmode asymptotics
    rigorously produce the factor~$16$ in
    $C^* = 16\cdot 4\pi^2/3$, and the $\lambda=11$ discriminator does not
    persist as a separate branch, then
    Theorem~GF1 (Part~II~\cite{Geiger2026partII}) provides a
    boundary-moment scaling benchmark for the non-exceptional Branch~A
    regime; the $\lambda=3$ branch remains a conservative guardrail.
\end{enumerate}

These two slots are \emph{independent}: closing one does not affect the
other. Each is a well-posed analytical problem amenable to standard tools
(complex analysis for B-1; Fourier analysis and spectral theory for GF1).

\begin{remark}[Both slots are margin-type --- a strategic caveat]
\label{rem:slots-margin}
Both slots are \emph{margin/coercivity} statements: a uniform bound that closes a
gap. They remain well-posed open problems, and nothing below refutes them.
\Cref{sec:sharpened-wall} adds a \emph{heuristic} criticality lens: since RH (if
true) sits at the critical value $\Lambda=0$ with no buffer
(\cref{sec:lambda-lens}), exact-structural mechanisms (identity, symmetry,
spectral realization) deserve \emph{additional} strategic weight alongside these
margin-type slots. This is a re-reading of priorities, not a theorem and not a
refutation (see the caveat in \cref{rem:lambda-caveat}).
\end{remark}


% =============================================================================
\section{The Sharpened Common Wall: A Criticality Lens and an Analytic-Continuation Localization}
\label{sec:sharpened-wall}
% =============================================================================

A 2026 working phase re-examined all branches through one question --- \emph{what
kind of argument can close the gap?} --- and localized the shared obstruction more
sharply. Everything in this section is a \textbf{strategic lens and a diagnostic
localization}, not a new closure of RH and not a refutation of any route. Several
items are honest self-corrections of internally proposed constructions. The
conditional-reduction status of Parts~II/III is unchanged.

\subsection{A criticality lens (heuristic, not a theorem)}
\label{sec:lambda-lens}

The de~Bruijn--Newman constant $\Lambda$ satisfies $\Lambda\le 0\Leftrightarrow$ RH;
Rodgers and Tao (\emph{Ann.\ of Math.}\ 2020) proved $\Lambda\ge 0$. Hence
$\mathrm{RH}\Leftrightarrow\Lambda=0$: if RH holds, it sits \emph{exactly} at the
critical value, with no margin to spare.

\begin{remark}[Criticality lens --- and its essential caveat]
\label{rem:lambda-caveat}
This invites a \emph{strategic} dichotomy of proof mechanisms:
\emph{margin-seeking} (a uniform gap / coercivity / smallness bound) versus
\emph{exact-structural} (an identity, a symmetry, or a spectral realization that
forces the conclusion without slack). The heuristic reading is that a mechanism
producing a strict uniform margin would be establishing something with a buffer,
which sits uneasily with a critical truth at $\Lambda=0$.
\textbf{Caveat (essential, to avoid over-claiming).} The step ``a uniform margin
would force $\Lambda<0$'' is an \emph{analogy}, not a theorem. The branch-specific
margins --- the even-dominance gap $\lambda_1^+<\lambda_1^-$, the Hadamard constant
$\kappa$, the CCM residual smallness --- are \emph{not} shown to equal a
de~Bruijn--Newman statement; $\Lambda$ and these quantities are different objects.
Consequently the lens does \emph{not} refute the well-posed proof slots of
\cref{sec:cross-route} (\cref{rem:slots-margin}); it only suggests weighting
exact-structural routes \emph{in addition}. It is a re-reading of priorities, not
a result.
\end{remark}

Under this lens the routes re-sort by \emph{type}, not by viability: the
margin-type branches~A (A7/AVG gap), B ((ii-a) coercivity), Z (residual
smallness) keep their exact kernels as assets, while the exact-structural lines
--- M (Bost--Connes\,$\leftrightarrow$\,Meyer, with the KMS phase transition at
$\beta=1$ playing the role of built-in criticality) and the global adelic trace
route (Connes--Consani scaling site / $\mathbb{F}_1$, where the only missing datum
is the Hodge-index positivity of an otherwise-available Lorentzian form) ---
receive additional strategic weight. The empirical signature is consistent: the
even-dominance margin grows like a power law $\sim\lambda^{0.72}$ (a criticality
signature, not a buffer). This re-weighting is recorded in the project decision
log; it does not alter any route's formal status.

\subsection{The common wall: analytic continuation across $\Re s=1$}
\label{sec:common-wall-ac}

Independently of the lens, the correctly-typed routes share a single concrete
endpoint. Three lines of work --- the M\"obius prime-hub, the Wiener--Hopf
factorization, and the Sonin/Toeplitz (Connes--Consani) construction --- are each
\emph{locally/finitely} controllable but break at \emph{exactly} the same place:
the global limit, where the local data must assemble into $\zeta$. The unified
description is the Weyl--Titchmarsh function $M_\infty$ of the unbounded limit
operator; the missing datum is uniform control of that limit from local data
\emph{without} importing $\zeta$.

\begin{remark}[The wall is a classical fact, here sharply localized]
\label{rem:wall-folklore}
A finite Euler product $\prod_{p\in P}(1-p^{-s})^{-1}$ has, after the substitution
$z=-\mathrm{i}(s-\tfrac12)$, a pole comb in the upper half-plane at
$z_{p,k}=2\pi k/\log p+\mathrm{i}/2$ (on $\Im z=\tfrac12$), with a common pole at
$z=\mathrm{i}/2$ of order $|P|$; hence its logarithmic derivative is not a Herglotz
function. The comb is removed only by $\zeta$'s analytic continuation
($\zeta(0)$ finite; the pole at $s=1$). This is \textbf{Euler-product folklore}
(the product converges only for $\Re s>1$). It is not a new obstruction: it is a
\emph{precise localization} of the known wall --- the missing datum is the
analytic continuation across $\Re s=1$, i.e.\ a global/functional-equation
phenomenon that finite arithmetic-plus-archimedean data structurally do not
supply.
\end{remark}

\subsection{A computable diagnostic from the classical Laguerre--P\'olya/Herglotz criterion}
\label{sec:diagnostic}

Write $\Xi(z)=\xi(\tfrac12+\mathrm{i}z)$ and $M_\Xi=-\Xi'/\Xi$. The chain
\[
\mathrm{RH}\ \Longleftrightarrow\ \Xi\in\text{Laguerre--P\'olya}
\ \Longleftrightarrow\ M_\Xi\ \text{Herglotz}
\ \Longleftrightarrow\ \tfrac{(M_\Xi(z)-\overline{M_\Xi(w)})}{z-\bar w}\succeq 0
\]
is \textbf{classical} (Hermite--Biehler; the Laguerre--P\'olya characterization of
real-rooted entire functions; the positivity $\Re\,\xi'/\xi>0$ for $\Re s>\tfrac12$
is Voros' criterion). \emph{We claim no novelty for the criterion itself.} The
contribution is a \emph{computable, non-circular test protocol} that turns this
classical equivalence into a diagnostic instrument: an Euler-anchor giving
$M_\Xi$ from $\xi'/\xi$ in the convergent region $\Re w>1$ (verified numerically to
$\sim 10^{-31}$), a ladder of gates (normalization; tail-tightness; a Stieltjes
identity; an off-axis Poisson tomography resolving individual residues), and a
detector lemma by which an off-line zero forces a negative $1\times1$ Pick minor
just below the migrated pole. The instrument is empirically validated: it detects
off-line zeros and discriminates the true zero spectrum from a smooth
density-matched control.

\begin{remark}[What the diagnostic established on a genuine operator]
An independent, non-circular discretization of the Connes--Moscovici prolate
operator (arXiv:2112.05500) was run through the full instrument. It reproduces
Riemann's zero-counting function $N(E)$ to $\sim 3\%$ in the UV --- the
\emph{archimedean half} (the smooth density, leading term $\tfrac12\log(y/2\pi)$),
an independent non-circular confirmation of the CCM density claim. But its poles
sit $O(1)$ \emph{off} the zeros, so the arithmetic fine-structure (the
\emph{positions}) is not captured; and the relative-determinant / Krein
construction does \emph{not} remove the Euler pole comb, because the prolate
background is itself Herglotz (comb-free), leaving nothing to cancel. Both findings
re-confirm \cref{rem:wall-folklore}: density yes, positions no; the wall is the
analytic continuation. The genuine positive by-product --- that the determinant
channel $-D'/D$ (residues~$=$~multiplicities), not mere eigenvalue convergence, is
the object that must converge ($\det_{\mathrm{reg}}(A-z)\to C\cdot\Xi$) --- sharpens
the CCM frontier and is recorded among the independent results
(\cref{sec:independent-results}).
\end{remark}

\subsection{Four closures of concrete constructions (not routes)}
\label{sec:four-closures}

The phase produced four genuine negative results, each over a \emph{specific
constructed object or technique}. None closes a whole route; each re-types the
route or names a successor.

\begin{enumerate}
  \item \textbf{Osterwalder--Schrader reflection, natural form.} The reflection
    block $Q_+(x)=\iint_0^\infty \overline{x(t)}\,K(t+s)\,x(s)$ is indefinite,
    because the prime distributions $\delta(t+s-\log p)$ in the Weil kernel meet
    the half-axis $t+s=\log p$; the Toeplitz alternative $K(t-s)$ collapses to
    global Weil positivity $=$ RH. The OS route in its most natural form
    (inversion reflection on $\{|g|\ge 1\}$) is closed-negative; another
    half-space choice would be needed.
  \item \textbf{M\"obius $\mathbb{Z}_2$ prime hub, na\"ive $L_s^-$.} An
    obstruction theorem: any continuation of the divergent prime traces
    $\operatorname{Tr}T_s=P(s)$, $\operatorname{Tr}T_s^2=P(2s)$ down to
    $\Re s=\tfrac12$ imports $\zeta$ and its zeros via
    $P(s)=\sum_n \mu(n)/n\cdot\log\zeta(ns)$ --- hence circular. (Project-internal;
    not a literature-verified standard.)
  \item \textbf{Toeplitz-symbol sub-method (Branch~A).} $A_{\sin}$ is not Toeplitz
    ($m$-dependent oscillating diagonals $\Rightarrow$ no well-defined symbol
    $f(\theta)$); a structural defect of the constructed object, not merely
    ``the method equals the open problem''. Successor: the finite Sylvester-inertia
    bound (done) plus GLT/outlier control.
  \item \textbf{Relative-determinant / Krein bypass.} The relative determinant
    $-\partial_z\log[\det(A-z)/\det(A^0-z)]$ does \emph{not} remove the Euler pole
    comb: if $A^0$ is the prolate background it is itself comb-free (nothing to
    cancel, and $\det(\text{Euler}\cdot\text{prolate})/\det(\text{prolate})
    =\det(\text{Euler})$ keeps the comb); the only comb-free alternative needs a
    self-adjoint prime factor, which is circular ($=$ RH). This is an honest
    self-correction of an internally proposed ``right technology''; the
    proposed-then-refuted arc is recorded in \cref{sec:sammler-disziplin}.
\end{enumerate}

\noindent
Crucially, the \emph{determinant channel} $-D'/D$ as a \emph{sharpening} of the CCM
frontier remains valid (\cref{sec:independent-results}); only the \emph{relative}
determinant as a \emph{bypass} of the continuation wall is refuted. The two must
not be conflated.


% =============================================================================
\section{Dormant Routes}
\label{sec:dormant}
% =============================================================================

A small number of constructed routes are dormant rather than dead: they
are explicit reductions or partial constructions that have not been
refuted, but currently lack active development. The internal register
(\texttt{AKTIONSPLAN.md}) keeps the full inventory; the four below are
summarised here because each has a direct bridge to one of the active
branches.

\begin{itemize}
  \item \textbf{D-CCM (Dirichlet--CCM operator
    $D^{(\lambda,N,\chi)}_{\log}$).} A direct $\chi$-twist of the
    CCM rank-$N$ perturbation construction. Addresses GRH, not RH
    directly --- the trivial character $\chi_0$ (Riemann case) is in
    this scaling regime not practically testable. Companion of
    Branch~Z.

  \item \textbf{D-Twist ($\chi$-twisted defect norm).} Diagnostic
    blueprint with completed Pilot~v2 and Milestone~2$_\chi$ work;
    reformulated as $\lambda$-asymptotic statement. First tests at
    $\chi_{21}, \chi_{33}$; Riemann unreachable in this scale.

  \item \textbf{D-PW (Paley--Wiener sector difference
    $\psi_\chi = \phi_\chi^+ - \phi_\chi^-$).} Classification of the
    five Paley--Wiener steps complete: S1, S2, S4 mechanical; S3, S5
    character-specific. Continuation of the published Dirichlet Character
    Atlas v2 (Zenodo~\texttt{20241612}).

  \item \textbf{P-M (M\"obius-graded $\mathbb{Z}_2$ prime hub).}
    Conditional lemma proved for finite prime set (M\"obius identity
    exact); full prime set converges only for $\Re(s)>1$;
    $\mathbb{Z}_2$-sector statement for non-trivial zeros open.
    Structurally appealing alternative interpretation of the Prime Orbit
    Axiom (Part~I, Definition~2.1). A 2026 obstruction theorem sharpens
    this (\cref{sec:four-closures}): no non-circular $L_s^-$ exists, since
    the divergent prime traces continue to $\Re s=\tfrac12$ only by
    importing $\zeta$ and its zeros. This re-types P-M as the \emph{prime
    half} of a P-M$+$M fusion --- it provably lacks the archimedean--modular
    half ($s\leftrightarrow 1-s$) that route~M supplies (project-internal).
\end{itemize}

The remaining dormant entries (P-A, P-C, P-B, M, R, H) are not listed
individually here; see \texttt{AKTIONSPLAN.md} for the full inventory. Two are
\emph{strategically elevated} by the criticality lens of
\cref{sec:lambda-lens}: \textbf{M} (Bost--Connes\,$\leftrightarrow$\,Meyer) and
the global adelic trace route \textbf{P-A} (Connes--Consani scaling site) are of
the exact-structural type and now receive additional strategic weight; \textbf{P-C}
is co-resolved by an active branch; \textbf{P-B} (Berry--Keating) is of the right
type but unconstructed (self-adjointness is the missing Hilbert--P\'olya theorem);
\textbf{R}, \textbf{H} are framework/expository.


% =============================================================================
\section{External Lines}
\label{sec:external}
% =============================================================================

The landscape sits inside a wider research ecosystem. The following lines
are cited as anchors; only the directly proof-relevant items appear here.

\begin{itemize}
  \item \textbf{Connes 2026 Mainline} \cite{Connes2026}
    (arXiv:2602.04022). Re-frames the Riemann zeros as the infrared
    spectrum of a self-adjoint operator and identifies the key open
    hypothesis MS2 (\S6.6(b)). Branches~A, B, Z all condition on this
    line.

  \item \textbf{Connes--van~Suijlekom 2025} \cite{ConnesvS2025}
    (arXiv:2511.23257). Theorem~5.6: 1-dim even kernel of the truncated
    Weil form implies real zeros of $\hat\eta$. This is A1~=~B1 in the
    present landscape.

  \item \textbf{Connes--Consani--Moscovici 2025}
    \cite{ConnesMoscovici2023} (arXiv:2511.22755). Foundation of the
    rank-one perturbation $D^{(\lambda,N)}_{\log}$. Anchors Branch~B's
    Foundation-Lock and Branch~Z's microcluster construction.

  \item \textbf{Groskin 2026} (arXiv:2605.20224). Independent
    implementation of the CCM Galerkin matrix at sixteen cutoffs
    $c=13\ldots 67, 100$ at high precision (Aitken extrapolation,
    $\sim 10^{-334}$ depth). Acts as external validator for branches~B
    and~Z.

  \item \textbf{\'Sliwi\'nski 2026} (arXiv:2601.12133). Counter-regime:
    in the $\lambda = N$ configuration with low precision, proves an
    inverse-logarithmic dissonance lower bound
    $\epsilon(\lambda,N)\ge 1/(4\ln\lambda)$. Does not bound the
    fixed-$\lambda$, $N\gg\lambda$, high-precision setup of this
    series; stands as a warning against na\"ive $\lambda = N$ limits.

  \item \textbf{Watson 2026} (arXiv:2602.01248), ``The Riemann
    $\Xi$-function from primitive Markovian cycles~I.'' Constructs a
    Laguerre--P\'olya function $\Psi$ from finite reversible Markov dynamics,
    ``entirely Archimedean, using no Euler products, primes, or arithmetic
    spectral input''; not an RH proof --- it isolates ``a single remaining
    analytic problem relating $\Psi$ to $\Xi(2\cdot)$.'' An independent
    triangulation of the common wall: from the archimedean side it reaches the
    same \emph{identification} step (relating the constructed object to $\Xi$)
    that the routes of this series reach from the arithmetic side
    (\cref{sec:common-wall-ac}).

  \item \textbf{Franchi Vicer\'e 2026} (Zenodo~19546495). Claims
    unconditional RH proof via apeirohedron / semi-local spectral
    descent. Surveyed and found not viable: the stabilisation argument
    confuses pointwise with strip-uniform convergence; the Apeirohedron
    limit step is not rigorous; boundary conditions do not replace the
    even-sector / MS1 input. See \texttt{AKTIONSPLAN.md} \S Klasse~8.
\end{itemize}


% =============================================================================
\section{Systematically Explored Alternatives}
\label{sec:explored-alternatives}
% =============================================================================

Eleven alternative strategies (BI-1 through BI-11) were systematically
explored during development. This exhaustive search is itself a
meta-result: the solution space is mapped.

\begin{center}
\footnotesize
\renewcommand{\arraystretch}{1.2}
\resizebox{\textwidth}{!}{%
\begin{tabular}{l|p{4.8cm}|l|l}
\textbf{Approach} & \textbf{Core idea} & \textbf{Result} & \textbf{Status} \\\hline
BI-1: $\det_2$ / Carleman & Extend to Hilbert--Schmidt class & No positivity substitute & Refuted \\
BI-2: RG fixed point & Weil positivity under scaling & Not formalizable & Discarded \\
BI-3: Universal property & Physical plausibility $\Rightarrow$ RH & Too abstract & Discarded \\
BI-5: Lifting programs & Weil $\to$ Grothendieck $\to$ RMT & All open except $\mathbb{F}_q$ & Not usable \\
BI-6: Implication ring & Li $\leftrightarrow$ Weil $\leftrightarrow$ NB $\leftrightarrow$ BD & Ring does not exist & Refuted \\
BI-7: P\"oschl--Teller scattering & Gamma structure analog to $\xi$ & No off-line stability & Negative \\
\textbf{BI-8: Second-order stability} & \textbf{RH as variational problem} & \textbf{Architecture correct} & \textbf{Realized} \\
BI-9: Soliton models & Topological stability & Coulomb gas unstable & Refuted \\
\textbf{BI-10: Building blocks} & \textbf{No-Coordination, shift structure} & \textbf{Successfully reused} & \textbf{Integrated} \\
\textbf{BI-11: Connes 2026} & \textbf{Even dominance via $Q_W$} & \textbf{Breakthrough} & \textbf{Core path}
\end{tabular}}
\end{center}

\noindent The chosen path (BI-8 + BI-10 + BI-11) combines the
variational architecture, the arithmetic building blocks, and Connes'
spectral framework. Routes~A and~B (total positivity, commuting operator)
were ruling-outs that are structural discoveries: even dominance holds
\emph{despite} their absence, which clarifies that the mechanism is
genuinely arithmetic-statistical rather than algebraic
(Part~II~\cite{Geiger2026partII}).


% =============================================================================
\section{Independent Results}
\label{sec:independent-results}
% =============================================================================

Several results from this program are of independent interest, regardless
of the main proof's status:

\begin{center}
\footnotesize
\renewcommand{\arraystretch}{1.2}
\resizebox{\textwidth}{!}{%
\begin{tabular}{p{4.2cm}|p{3.2cm}|p{4.2cm}|l}
\textbf{Result} & \textbf{Context} & \textbf{Significance} & \textbf{Status} \\\hline
\multicolumn{4}{l}{\emph{Positive results}} \\[2pt]
Shift Parity Lemma & Arithmetic trigonometry & Every prime favors even eigenfunctions & Proved \\
No-Coordination Lemma & Multiplicative independence & Arithmetic substitute for entropy term & Proved \\
Leading-Mode Cancellation ($c = 2+\sqrt{2}$) & Resolvent analysis & Universal pairwise cancellation constant & Proved \\
Euler--Maclaurin: $\rho^{\mathrm{EM}} < 0$ & PNT asymptotics & Global stability anchor & Proved \\
$33$ finite-range CAP diagnostics & Interval arithmetic & Conditional bridge & Evidence \\
Concentration lemma & Even-dominance image measure & No fixed prime set controls a positive share of prime mass & Proved \\
Finite Sylvester inertia ($\nu_-=0$) & Interval arithmetic, $\lambda=100$ & Finite even dominance rigorous, no symbol needed & Verified \\
Residue-calibration channel & Weyl-function diagnostic & Eigenvalue conv.\ $\neq$ Weyl conv.; channel is $-D'/D$, $\det_{\mathrm{reg}}\!\to C\Xi$ & Diagnostic \\[4pt]
\multicolumn{4}{l}{\emph{Negative results}} \\[2pt]
Obstructions K1--K4 & Dead-end analysis & Maps impossible strategies & Documented \\
No absolute $\mathrm{PF}_\infty$ for $A_\lambda$ & Prime shift multiplier & $M(\xi)$ negative on $\sim49\%$ of domain & Proved \\
No universal commuting operator & Sturm--Liouville search & Kernel of $[T, D_N(r)]$ is $\{c \cdot I\}$ & Proved \\[4pt]
\multicolumn{4}{l}{\emph{Structural insights}} \\[2pt]
Inverted-Hessian structure & YM/TU/RH comparison & RH is a maximum, not minimum & Confirmed \\
Cross-disciplinary catalog & Regularization & $11$ fields share trace-class obstruction & Documented
\end{tabular}}
\end{center}

\noindent The negative results protect future researchers from repeating
dead-end strategies. The Shift Parity Lemma, the obstruction analysis
(K1--K4), and the Mellin decay barrier are publishable independently.


% =============================================================================
\section{Hypothesis Corrections in the Iterative Process}
\label{sec:sammler-disziplin}
% =============================================================================

The iterative development produced several hypothesis corrections that
are documented here as a record of honest research practice
(Sammler-Disziplin).

\subsection{GF1 Correction Chain (Iterations 15--29)}

\begin{center}
\renewcommand{\arraystretch}{1.2}
\resizebox{\textwidth}{!}{%
\begin{tabular}{c|l|l}
\toprule
\textbf{Iter} & \textbf{Hypothesis} & \textbf{Status} \\
\midrule
15 & $C_m \approx 8$ universal & falsified (boundary effect) \\
18 & $C_m(\lambda=5) \approx 4.3$ & falsified (no plateau) \\
19 & $\mathrm{sLI} \sim 1/N$ & confirmed, but without $\lambda$-scaling \\
20 & $N \cdot \mathrm{sLI} \approx 329$ universal & falsified \\
\textbf{27} & $\mathbf{\mathrm{sLI} = C^*\,\lambda/(NL) + o(\cdot)}$ & confirmed for $\lambda=5,7$, $C^* = 64\pi^2/3$ \\
28 & $\lambda=3$ all-$\lambda$ check & low-bandwidth outlier; two-regime guardrail \\
29 & $\lambda=11$ discriminator and endpoint-kernel audit & $C_{\mathrm{eff}}\approx168$ at $N=200$; simple endpoint factorisation falsified \\
\bottomrule
\end{tabular}
}
\end{center}


\subsection{\texorpdfstring{B-1 Corrections and $c_\Xi$ RH-Neutral Form}{B-1 Corrections and cXi RH-Neutral Form}}

The Hadamard constant $c_\Xi$ was initially written in the RH-specialised
form $\sum_{k\geq 1} 1/\gamma_k^2$ (sum over imaginary parts of zeros on
the critical line). An external audit corrected this to the RH-neutral
form
$c_\Xi = -\tfrac{1}{2}\sum_\rho 1/(\rho-\tfrac{1}{2})^2$
(sum over all non-trivial zeros). This is a Sammler-Disziplin correction:
the RH-neutral form is mathematically more general and does not assume what
is being proved.


\subsection{The Relative-Determinant Bypass: Proposed and Refuted (2026)}

During the 2026 phase, the relative determinant (Krein spectral shift)
$-\partial_z\log[\det(A-z)/\det(A^0-z)]$ was proposed internally as the ``right
technology'' to remove the upper-half-plane Euler pole comb and so bypass the
continuation wall (\cref{rem:wall-folklore}). A subsequent test on the genuine
Connes--Moscovici prolate background refuted this as a bypass: the prolate $A^0$ is
itself Herglotz (comb-free), so there is nothing to cancel
($\det(\text{Euler}\cdot\text{prolate})/\det(\text{prolate})=\det(\text{Euler})$
retains the comb), and the only comb-free alternative needs a self-adjoint prime
factor, which is circular ($=$~RH). The proposal is withdrawn as a bypass. The
\emph{separate} finding --- that the determinant channel $-D'/D$, not mere
eigenvalue convergence, is the correct object to converge
($\det_{\mathrm{reg}}(A-z)\to C\,\Xi$) --- remains valid
(\cref{sec:independent-results}); the two must not be conflated.

A second correction: an internal phase lemma
$\Theta_\sigma=\xi(\sigma-\mathrm{i}z)/\xi(\sigma+\mathrm{i}z)$ was sign-reversed
(its poles come from the denominator, so under RH it \emph{has} poles); the correct
object is the reciprocal phase
$\widehat\Theta_\sigma=\xi(\sigma+\mathrm{i}z)/\xi(\sigma-\mathrm{i}z)$, with
RH~$\Leftrightarrow\widehat\Theta_\sigma$ pole-free in the upper half-plane for all
$\sigma>\tfrac12$.


\subsection{Value of the Correction Record}

These corrections demonstrate that the research process is
self-correcting. Each correction sharpened the theoretical framework:
the GF1 chain moved from a na\"ive constant ($C_m \approx 8$) through
several intermediate forms to the $\lambda=5,7$ boundary-moment
benchmark $\mathrm{sLI} \sim C^*\lambda/(NL)$; the $\lambda=3$ sweep and
the open $\lambda=11$ discriminator prevent an all-$\lambda$ theorem reading,
and the $c_\Xi$ correction removed a
circular RH-dependence from a key constant.


% =============================================================================
\section{Update (v3.2): The RH/abc Archimedean-Mass Crosswalk}
\label{sec:v32-abc}
% =============================================================================

A cross-problem line that had been carried informally alongside the routes
of this series was audited to a close between June and August 2026. It is
recorded here in Part~IV rather than with the spectral material of
Part~III~\cite{Geiger2026partIII}, because it concerns a bridge between RH
and a second conjecture rather than a route inside the RH landscape
itself.

The crosswalk (register designation~B2, state~V1.1) compared the
archimedean mass objects on the two sides: the RH-side $\mu_-$, a signed
dimensionless negative mass, against the abc-side $\Omega_E$, a positive
N\'eron period. Only a structural analogy is confirmed --- the two are not
instances of a common object, and no shared normalisation emerged under
which they could be compared quantitatively. An accompanying candidate
audit examined four proposals for such a common object, namely the
Connes--Consani trace, the Yuan--Zhang height pairing, the Frey-period
log-height, and a formal container, and found that none of them satisfies
the four bridge obligations. The audit rests on five bound sources with
producer and verifier quotas of $19/19$ each, fail-closed. The consequence
for the landscape is that the RH\,$\leftrightarrow$\,abc bridge is
exhausted \emph{as a linguistic analogy}: a successor would have to supply
a normalised common $H_\infty$ with typed support, both realisations,
compatible normalisations, and two non-circular inequalities. Claim
boundary: what closes here is the candidate list examined and the analogy
reading of the bridge --- neither conjecture is refuted, and an object
outside the four tested candidates is untested rather than excluded.


% =============================================================================
\section{Summary}
\label{sec:summary}
% =============================================================================

Part~IV contributes the cross-route perspective: the three active
branches share a Galerkin pattern (fixed-parameter positive, uniform-limit
open), and two isolated proof slots (B-1 log-curvature, GF1
factor-16 identification) are the concrete analytical targets, which remain
well-posed open problems. \Cref{sec:sharpened-wall} then adds a 2026
\emph{heuristic} criticality lens ($\mathrm{RH}\Leftrightarrow\Lambda=0$) that
re-weights the routes toward exact-structural mechanisms (without refuting the
margin-type slots), localizes the common wall as the analytic continuation across
$\Re s=1$, records a computable Laguerre--P\'olya/Herglotz diagnostic (a
classical criterion turned into a test instrument), and reports four closures of
\emph{concrete constructions} --- the OS reflection in natural form, the na\"ive
M\"obius $L_s^-$, the Toeplitz-symbol sub-method, and the relative-determinant
bypass --- none of which closes a route or RH. The dormant
and external routes are catalogued for completeness; the hypothesis
correction record provides transparency. The condensed atlas and final
assessment follow in Part~V~\cite{Geiger2026partV}.


\bigskip
\noindent\textbf{Acknowledgments.} Computations were performed on a
Mac~Studio (Apple M1~Max, 32~GB) and a Hetzner CCX13 cloud server
(2~dedicated vCPU, 8~GB). The interval arithmetic library mpmath was
developed by Fredrik Johansson.


\section*{AI Disclosure}

This manuscript was developed in extensive collaboration with AI language
models (Claude, Anthropic; Copilot, Microsoft/GitHub; Gemini, Google
DeepMind). AI contributed to conceptual exploration, analytical work,
literature review, writing, and critical review. The author, Lukas
Geiger, conceived the research direction, provided domain expertise, and
exercised final editorial authority. The author assumes full and sole
scholarly responsibility for all content, including any errors.


\bibliographystyle{plain}
\bibliography{fst-rh-references}

\end{document}
